\documentclass[12pt]{article}
\usepackage{amsmath,amssymb}

\begin{document}
\section*{{2020 AMC 12A Problems/Problem 12}}
Source: AoPS Wiki

\subsection*{Solution}
The slope of the line is $\frac{3}{5}$ . We must transform it by $45^{\circ}$ . $45^{\circ}$ creates an isosceles right triangle, since the sum of the angles of the triangle must be $180^{\circ}$ and one angle is $90^{\circ}$ . This means the last leg angle must also be $45^{\circ}$ . In the isosceles right triangle, the two legs are congruent. We can therefore construct an isosceles right triangle with a line of $\frac{3}{5}$ slope on graph paper. That line with $\frac{3}{5}$ slope starts at $(0,0)$ and will go to $(5,3)$ , the vector $<5, 3 >$ . Construct another line from $(0,0)$ to $(3,-5)$ , the vector $<3,-5>$ . This is $\perp$ and equal to the original line segment. The difference between the two vectors is $<2,8>$ , which is the slope $4$ , and that is the slope of line $k$ . Furthermore, the equation $3x-5y+40=0$ passes straight through $(20,20)$ since $3(20)-5(20)+40=60-100+40=0$ , which means that any rotations about $(20,20)$ would contain $(20,20)$ . We can create a line of slope $4$ through $(20,20)$ . The $x$ -intercept is therefore $20-\frac{20}{4} = \boxed{\textbf{(B) } 15.}$ ~lopkiloinm ~ShawnX (diagram)

\end{document}
